% \begin{document}

\chapter{Denotational Semantics}

We have now seen two operational models for programming languages: small-step and large-step.
In this lecture, we consider a different semantic model, called \emph{denotational semantics}.

The idea in denotational semantics is to express the meaning of a program as the mathematical
function that expresses what the program computes. We can think of an IMP program $c$ as a
function from stores to stores: given an initial store, the program produces a final store. For
example, the program $\mathit{foo} := \mathit{bar} + 1$ can be thought of as a function that
when given an input store $\sigma$, produces a final store $\sigma'$ that is identical to
$\sigma$ except that it maps $\mathit{foo}$ to the integer $\sigma(\mathit{bar}) + 1$; that is,
$\sigma' = \sigma[\mathit{foo} \mapsto \sigma(\mathit{bar}) + 1]$. We will model programs as
functions from input stores to output stores. As opposed to operational models, which tell us
\emph{how} programs execute, the denotational model shows us \emph{what} programs compute.

\section{A Denotational Semantics for IMP}

For each program $c$, we write $\mathcal{C}[\![c]\!]$ for the \emph{denotation} of $c$, that
is, the mathematical function that $c$ represents:
\[
\mathcal{C}[\![c]\!] : \mathbf{Store} \rightharpoonup \mathbf{Store}.
\]
Note that $\mathcal{C}[\![c]\!]$ is actually a partial function (as opposed to a total
function), both because the store may not be defined on the free variables of the program and
because the program may not terminate for certain input stores. The function
$\mathcal{C}[\![c]\!]$ is not defined for non-terminating programs as they have no
corresponding output stores.

We will write $\mathcal{C}[\![c]\!]\sigma$ for the result of applying the function
$\mathcal{C}[\![c]\!]$ to the store $\sigma$. That is, if $f$ is the function that
$\mathcal{C}[\![c]\!]$ denotes, then we write $\mathcal{C}[\![c]\!]\sigma$ to mean the same
thing as $f(\sigma)$.

We must also model expressions as functions, this time from stores to the values they
represent. We will write $\mathcal{A}[\![a]\!]$ for the denotation of arithmetic expression
$a$, and $\mathcal{B}[\![b]\!]$ for the denotation of boolean expression $b$.
\[
\mathcal{A}[\![a]\!] : \mathbf{Store} \rightharpoonup \mathbf{Int}
\qquad
\mathcal{B}[\![b]\!] : \mathbf{Store} \rightharpoonup \{\mathbf{true}, \mathbf{false}\}
\]

Now we want to define these functions. To make it easier to write down these definitions, we
will describe (partial) functions using sets of pairs. More precisely, we will represent a
partial map $f : A \rightharpoonup B$ as a set of pairs $F = \{(a, b) \mid a \in A \text{ and }
b = f(a) \in B\}$ such that, for each $a \in A$, there is at most one pair of the form $(a,
\underline{\ \ })$ in the set. Hence $(a, b) \in F$ is the same as $b = f(a)$.

We can now define denotations for IMP. We start with the denotations of expressions:
\begin{align*}
\mathcal{A}[\![n]\!] &= \{(\sigma, n)\} \\
\mathcal{A}[\![x]\!] &= \{(\sigma, \sigma(x))\} \\
\mathcal{A}[\![a_1 + a_2]\!] &= \{(\sigma, n) \mid (\sigma, n_1) \in \mathcal{A}[\![a_1]\!]
\wedge (\sigma, n_2) \in \mathcal{A}[\![a_2]\!] \wedge n = n_1 + n_2\}
\end{align*}
\begin{align*}
\mathcal{B}[\![\mathbf{true}]\!] &= \{(\sigma, \mathbf{true})\} \\
\mathcal{B}[\![\mathbf{false}]\!] &= \{(\sigma, \mathbf{false})\} \\
\mathcal{B}[\![a_1 < a_2]\!] &= \{(\sigma, \mathbf{true}) \mid (\sigma, n_1) \in
\mathcal{A}[\![a_1]\!] \wedge (\sigma, n_2) \in \mathcal{A}[\![a_2]\!] \wedge n_1 < n_2\} \\
&\quad \cup \{(\sigma, \mathbf{false}) \mid (\sigma, n_1) \in \mathcal{A}[\![a_1]\!] \wedge
(\sigma, n_2) \in \mathcal{A}[\![a_2]\!] \wedge n_1 \geq n_2\}
\end{align*}

The denotations for commands are as follows:
\begin{align*}
\mathcal{C}[\![\mathbf{skip}]\!] &= \{(\sigma, \sigma)\} \\
\mathcal{C}[\![x := a]\!] &= \{(\sigma, \sigma[x \mapsto n]) \mid (\sigma, n) \in
\mathcal{A}[\![a]\!]\} \\
\mathcal{C}[\![c_1; c_2]\!] &= \{(\sigma, \sigma') \mid \exists \sigma''.\ ((\sigma, \sigma'')
\in \mathcal{C}[\![c_1]\!] \wedge (\sigma'', \sigma') \in \mathcal{C}[\![c_2]\!])\}
\end{align*}

All these can be expressed in function-style notation too.\\

\nt{
  If the denotation of $ \sigma $ under $ b $ is undefined, then the function should also be undefined, hence the case completeness.
}

Note that $\mathcal{C}[\![c_1; c_2]\!] = \mathcal{C}[\![c_2]\!] \circ \mathcal{C}[\![c_1]\!]$,
where $\circ$ is the composition of relations, defined as follows: if $R_1 \subseteq A \times
B$ and $R_2 \subseteq B \times C$ then $R_2 \circ R_1 \subseteq A \times C$ is $R_2 \circ R_1 =
\{(a, c) \mid \exists b \in B.\ (a, b) \in R_1 \wedge (b, c) \in R_2\}$. If
$\mathcal{C}[\![c_1]\!]$ and $\mathcal{C}[\![c_2]\!]$ are total functions, then $\circ$ is
function composition.

\begin{align*}
\mathcal{C}[\![\mathbf{if}\ b\ \mathbf{then}\ c_1\ \mathbf{else}\ c_2]\!] &= \{(\sigma,
\sigma') \mid (\sigma, \mathbf{true}) \in \mathcal{B}[\![b]\!] \wedge (\sigma, \sigma') \in
\mathcal{C}[\![c_1]\!]\} \\
&\quad \cup \{(\sigma, \sigma') \mid (\sigma, \mathbf{false}) \in \mathcal{B}[\![b]\!] \wedge
(\sigma, \sigma') \in \mathcal{C}[\![c_2]\!]\}
\end{align*}

\begin{align*}
\mathcal{C}[\![\mathbf{while}\ b\ \mathbf{do}\ c]\!] &= \{(\sigma, \sigma) \mid (\sigma,
\mathbf{false}) \in \mathcal{B}[\![b]\!]\} \\
&\quad \cup \{(\sigma, \sigma') \mid (\sigma, \mathbf{true}) \in \mathcal{B}[\![b]\!] \wedge
\exists \sigma''.\ ((\sigma, \sigma'') \in \mathcal{C}[\![c]\!] \wedge (\sigma'', \sigma') \in
\mathcal{C}[\![\mathbf{while}\ b\ \mathbf{do}\ c]\!])\}
\end{align*}

But now we have a problem: the last ``definition'' is not really a definition because it
expresses $\mathcal{C}[\![\mathbf{while}\ b\ \mathbf{do}\ c]\!]$ in terms of itself! This is not
a definition but a recursive equation. What we want is the solution to this equation.

\nt{
  The second invocation of while isn't completely circular because the state changes.
}

\section{Fixed Points}

We gave a recursive equation that the function $\mathcal{C}[\![\mathbf{while}\ b\ \mathbf{do}\
c]\!]$ must satisfy. To understand some of the issues involved, let's consider a simpler
example. Consider the following equation for a function $f : \mathbb{N} \to \mathbb{N}$.
\begin{equation}
f(x) =
\begin{cases}
0 & \text{if } x = 0 \\
f(x - 1) + 2x - 1 & \text{otherwise}
\end{cases}
\end{equation}
This is not a definition for $f$, but rather an equation that we want $f$ to satisfy. What
function, or functions, satisfy this equation for $f$? The only solution to this equation is
the function $f(x) = x^2$.

In general, there may be no solutions for a recursive equation (e.g., there are no functions $g
: \mathbb{N} \to \mathbb{N}$ that satisfy the recursive equation $g(x) = g(x) + 1$), or multiple
solutions (e.g., find two functions $g : \mathbb{R} \to \mathbb{R}$ that satisfy $g(x) = 4
\times g(\frac{x}{2})$).

We can compute solutions to such equations by building successive approximations. Each
approximation is closer and closer to the solution. To solve the recursive equation for $f$, we
start with the partial function $f_0 = \emptyset$ (i.e., $f_0$ is the empty relation; it is a
partial function with the empty set for its domain). We compute successive approximations
using the recursive equation.
\begin{align*}
f_0 &= \emptyset \\
f_1 &=
\begin{cases}
0 & \text{if } x = 0 \\
f_0(x - 1) + 2x - 1 & \text{otherwise}
\end{cases}
\; = \{(0, 0)\} \\
f_2 &=
\begin{cases}
0 & \text{if } x = 0 \\
f_1(x - 1) + 2x - 1 & \text{otherwise}
\end{cases}
\; = \{(0, 0), (1, 1)\} \\
f_3 &=
\begin{cases}
0 & \text{if } x = 0 \\
f_2(x - 1) + 2x - 1 & \text{otherwise}
\end{cases}
\; = \{(0, 0), (1, 1), (2, 4)\} \\
&\ \ \vdots
\end{align*}
This sequence of successive approximations $f_i$ gradually builds the function $f(x) = x^2$.

We can model this process of successive approximations using a higher-order function $F$ that
takes one approximation $f_k$ and returns the next approximation $f_{k+1}$:
\[
F : (\mathbb{N} \rightharpoonup \mathbb{N}) \to (\mathbb{N} \rightharpoonup \mathbb{N})
\]
where
\[
(F(f))(x) =
\begin{cases}
0 & \text{if } x = 0 \\
f(x - 1) + 2x - 1 & \text{otherwise}
\end{cases}
\]

A solution to the recursive equation~(1) is a function $f$ such that $f = F(f)$. In general,
given a function $F : A \to A$, we have that $a \in A$ is a \emph{fixed point} of $F$ if $F(a) =
a$. We also write $a = \operatorname{fix}(F)$ to indicate that $a$ is a least fixed point of
$F$.

So the solution to the recursive equation~(1) is a fixed-point of the higher-order function
$F$. We can compute this fixed point iteratively, starting with $f_0 = \emptyset$ and at each
iteration computing $f_{k+1} = F(f_k)$. The fixed point is the limit of this process:
\begin{align*}
f &= \operatorname{fix}(F) \\
&= f_0 \cup f_1 \cup f_2 \cup f_3 \cup \ldots \\
&= \emptyset \cup F(\emptyset) \cup F(F(\emptyset)) \cup F(F(F(\emptyset))) \cup \ldots \\
&= \bigcup_{i \geq 0} F^i(\emptyset)
\end{align*}

\section{Denotation for Loops}

We can now write the correct denotation case for while loops as the fixed point of a
higher-order function:
\[
\mathcal{C}[\![\mathbf{while}\ b\ \mathbf{do}\ c]\!] = \operatorname{fix}(F)
\]
where
\[
F(f) = \{(\sigma, \sigma) \mid (\sigma, \mathbf{false}) \in \mathcal{B}[\![b]\!]\} \cup
\{(\sigma, \sigma') \mid (\sigma, \mathbf{true}) \in \mathcal{B}[\![b]\!] \wedge \exists
\sigma''.\ (\sigma, \sigma'') \in \mathcal{C}[\![c]\!] \wedge (\sigma'', \sigma') \in f\}
\]

\section{Kleene's Fixed Point Theorem}

Next we'll prove Kleene's fixpoint theorem, which shows that the fixed point used to define the
semantics of while commands exists.

\dfn{Scott Continuity}{
A function $F$ from $U$ to $U$ is said to be \emph{Scott-continuous} if for every chain $X_1
\subseteq X_2 \subseteq \ldots$ we have $F(\bigcup_i X_i) = \bigcup_i F(X_i)$.
}

It is not hard to show that if $F$ is Scott continuous, then it is also \emph{monotonic}---that
is, $X \subseteq Y$ implies $F(X) \subseteq F(Y)$. The proof of this fact is left as an
exercise.

\thm{Kleene Fixpoint}{
Let $F$ be a Scott-continuous function. The least fixed point of $F$ is $\bigcup_i F^i
(\emptyset)$.
}

\pf{}{
Let $X = \bigcup_i F^i(\emptyset)$.

First, we will prove that $X$ is a fixed point of $F$---that is, $F(X) = X$. We calculate as
follows:
\begin{align*}
F(X) &= F\left(\bigcup_i F^i(\emptyset)\right) & \text{By definition of } X \\
&= \bigcup_i F(F^i(\emptyset)) & \text{By Scott continuity} \\
&= \bigcup_i F^{i+1}(\emptyset) \\
&= \emptyset \cup \bigcup_i F^{i+1}(\emptyset) \\
&= F^0(\emptyset) \cup \bigcup_i F^{i+1}(\emptyset) \\
&= \bigcup_i F^i(\emptyset) \\
&= X
\end{align*}

Second, we will show that $X$ is the least fixed point of $F$. Suppose that $Y$ is some other
arbitrary fixed point of $F$.

By induction, we can easily show that $F^i(\emptyset) \subseteq Y$ for all $i$. For the base
case, $i$ is $0$ and we trivially have $F^0(\emptyset) = \emptyset \subseteq Y$. For the
inductive case, we assume that $F^i(\emptyset) \subseteq Y$ and prove that $F^{i+1}(\emptyset)
\subseteq Y$. By our inductive hypothesis and the fact that $F$ is monotone, we have that
$F(F^i(\emptyset)) \subseteq F(Y)$. As $Y$ is a fixed point we also have $F(Y) = Y$ and so
$F^{i+1}(\emptyset) \subseteq Y$.

Then, since every element of the chain
\[
F^0(\emptyset) \subseteq F^1(\emptyset) \subseteq \ldots
\]
is a subset of $Y$, immediately we have that their union, $X = \bigcup_i F^i(\emptyset)
\subseteq Y$. Hence, $X$ is the least (with respect to $\subseteq$) fixed point of $F$.
}

\section{Reasoning using Denotational Semantics}

One of the main advantages of using denotational semantics compared to operational semantics
is that proofs of equivalence can be carried out directly by simply calculating the denotations
of programs and then arguing that they are identical. This is in contrast to operational
techniques, where one must reason explicitly about low-level transitions and derivations
involving ad hoc abstract machines.

As an example, to show that $\mathbf{skip}; c$ and $c; \mathbf{skip}$ are equivalent, we can
calculate as follows,
\begin{align*}
\mathcal{C}[\![\mathbf{skip}; c]\!] &= \{(\sigma, \sigma'') \mid \exists \sigma'.\ (\sigma,
\sigma') \in \mathcal{C}[\![\mathbf{skip}]\!] \wedge (\sigma', \sigma'') \in \mathcal{C}[\![c]\!]\} \\
&= \{(\sigma, \sigma'') \mid (\sigma, \sigma'') \in \mathcal{C}[\![c]\!]\} \\
&= \{(\sigma, \sigma'') \mid \exists \sigma'.\ (\sigma, \sigma') \in \mathcal{C}[\![c]\!] \wedge
(\sigma', \sigma'') \in \mathcal{C}[\![\mathbf{skip}]\!]\} \\
&= \mathcal{C}[\![c; \mathbf{skip}]\!]
\end{align*}
using standard facts about partial functions, relations, and sets as convenient in the proof
itself.

Next, consider the command $\mathcal{C}[\![\mathbf{while}\ \mathbf{false}\ \mathbf{do}\ c]\!]$.
By the definition of the denotational semantics, this is equal to $\operatorname{fix}(F)$ where
\[
F(f) = \{(\sigma, \sigma) \mid (\sigma, \mathbf{false}) \in \mathcal{B}[\![b]\!]\} \cup
\{(\sigma, \sigma'') \mid (\sigma, \mathbf{true}) \in \mathcal{B}[\![b]\!] \wedge (\sigma,
\sigma') \in \mathcal{C}[\![c]\!] \wedge (\sigma', \sigma'') \in f\}
\]
By the Kleene fixpoint theorem we have that $\operatorname{fix}(F) = \bigcup_i F^i(\emptyset)$.
It is straightforward to show by induction that since the guard is false, for all $i$ we have
$F^i(\emptyset) = \{(\sigma, \sigma)\}$. It follows that $\operatorname{fix}(F) = \{(\sigma,
\sigma)\}$, which is just $\mathcal{C}[\![\mathbf{skip}]\!]$.

As an exercise, calculate $\mathcal{C}[\![\mathbf{while}\ \mathbf{true}\ \mathbf{do}\
\mathbf{skip}]\!]$ using the same technique.

% \end{document}
